\section{Conclusion}\label{sec:conclusion}

We have shown that every minimizer of a joint upper-tail problem with
Erdős–Rényi reference probability $p$ and constraints $t(H_j,W)\ge r_j^{m_j}$ for
$1\le j\le s$ is multipodal with ordered blocks. For entropy maximization
with fixed densities $e(W)=\eps$ and $t(H,W)=\tau$, the same conclusion
holds throughout the open oversaturated range
$\eps^m<\tau<\tau_{\max}(\eps)$, provided $\Delta(H)\ge2$.
Vertex insertions connect the two problems to global minimization of
$\mathcal F_W$ and $\mathcal F_{W,\eps,\tau}$, respectively; the pointwise ordering
and finiteness of their minimizers yield the ordered block structure.
For fixed densities, we introduce a constraint
qualification that establishes multipodality with ordered
blocks.  We then prove that the condition holds for optimizers
at every subgraph density in the oversaturated range.

\paragraph{Block structure and phase transitions.}  Our finiteness argument does not
control the number of blocks in the multipodal optimizers.
Previous work has established bounds on the number of blocks for specific
constrained graphs and parameter regions~\cite{ImKimLimNamYangBipodal,
KenyonRadinRenSadun2017,kenyon2014entropy,NeemanRadinSadun2023,NeemanRadinSadun2024,
RadinSadunEdge2Star,RadinSadunNearExtreme}.
A next step is to establish such bounds in more general settings.
For example, even for triangle upper tails, it remains open whether
two blocks suffice throughout the symmetry-breaking region.

For $d$-regular graphs with $d\ge2$, Lubetzky and
Zhao~\cite{lubetzky2015large} determined the boundary between constant and
nonconstant optimizers in the $(p,r)$-plane.
Determining this boundary for general $H$, and locating transitions
between different block counts would refine
our structural result. Even when the block count stays fixed, numerical
experiments for triangle upper tails suggest jumps in the block
parameters~\cite{ImKimLimNamYangBipodal}. Proving such transitions and
describing how the optimizing block parameters vary within and between
phases remain further questions.

\paragraph{Several fixed subgraph densities.}  The multipodality conclusion
for joint upper tails cannot extend to arbitrary finite families of
prescribed subgraph densities: the universality result of Cooper,
Kr\'al', and Martins~\cite{CooperKralMartins2018} allows such constraints
to force a non-multipodal entropy maximizer.
Our theorem for one fixed $H$-density in the open oversaturated range
leaves a natural question between these two settings. Which restrictions
on the number or form of the constrained graphs ensure multipodality,
and which allow non-multipodal optimizers?

\paragraph{Lower tails.}  A related question is whether every minimizer of
\[
  \min\bigl\{I_p(W):W\in\W_0,\ t(H,W)\le r^m\bigr\},
  \qquad 0<r<p<1,
\]
or its joint version with several subgraph densities, is multipodal.
For triangles, the limiting problem as $p\to0$ with $r/p$ fixed has only
constant or symmetric bipodal minimizers~\cite{HuangLi2026LowerTail}.
With fixed edge and triangle densities, symmetric bipodal phases also
occur in parts of the undersaturated region $\tau<\eps^3$ in the
$(\eps,\tau)$-plane~\cite{NeemanRadinSadun2024}.
In these phases, connections are stronger between the blocks than within
them, so the two row types cross. Establishing multipodality for general
lower-tail constraints at fixed $p$, or throughout the fixed-density
undersaturated region, would therefore require an argument that allows
crossing rows.
